\section*{الفصل \ref{ch:g11:antibiotic-resistance} --- البكتيريا ومقاومة المضادات الحيوية}

\begin{solution}{exo:g11:antibiotic-resistance:1}
\omterm{def:g11:mutations:mutation}{الطفرات} التي تنشأ أثناء انقسامات النسيلة، \omterm{prop:g11:antibiotic-resistance:variation}{والنقل الأفقي للمورّثات} ---
أي \omterm{def:g10:universal-dna:information}{DNA}، وهو بلازميد عادةً، يتلقى من بكتيريا أخرى.
\end{solution}

\begin{solution}{exo:g11:antibiotic-resistance:2}
جزيء يقتل البكتيريا أو يوقف نموها بتأثيره في بنية توجد عند البكتيريا
وتفتقر إليها \omterm{def:g10:cells-common-unit:cell}{الخلايا} البشرية (الجدار، \omterm{def:g10:cells-common-unit:organelle}{والريبوزوم} الجرثومي، \omterm{def:g11:enzymes-and-phenotype:enzyme}{وإنزيمات}
التضاعف الجرثومية). والزكام فيروسي، والفيروسات لا هدف من هذه الأهداف
عندها.
\end{solution}

\begin{solution}{exo:g11:antibiotic-resistance:3}
\omterm{def:g11:enzymes-and-phenotype:enzyme}{إنزيم} يهدم \omterm{def:g11:antibiotic-resistance:antibiotic}{المضاد الحيوي}؛ أو هدف طافر لا يعود الدواء يرتبط به؛ أو مضخة
تطرد الدواء (أو جدار يمنعه من الدخول).
\end{solution}

\begin{solution}{exo:g11:antibiotic-resistance:4}
حساسة (20 فوق 18)، لكن بالكاد. والأفضل D، وهامشه فوق عتبته الأوسع، ثم
A.
\end{solution}

\begin{solution}{exo:g11:antibiotic-resistance:5}
لا. فعلى أطباق النسخ تظهر المستعمرات المقاومة في المواضع نفسها على كل
طبق مضاد، \omterm{def:g10:cells-common-unit:cell}{فالخلايا} المقاومة كانت موجودة في المستعمرات الأصلية، وهي لم
تلق الدواء قط. \omterm{def:g11:antibiotic-resistance:antibiotic}{فالمضاد الحيوي} يكشف وينتقي؛ ولا يستحث.
\end{solution}

\begin{solution}{exo:g11:antibiotic-resistance:6}
30 انقسامًا: أي $2^{30} \approx 10^9$ \omterm{def:g10:cells-common-unit:cell}{خلية}. وقد أنتجها نحو $10^9$
انقسام: أي نحو 10 \omterm{def:g10:cells-common-unit:cell}{خلايا} مقاومة وسطيًا --- وأكثر إن جاءت \omterm{def:g11:mutations:mutation}{الطفرة} باكرًا.
\end{solution}

\begin{solution}{exo:g11:antibiotic-resistance:7}
حول اليوم الثالث، حين تكون \omterm{def:g10:cells-common-unit:cell}{الخلايا} الحساسة قد هبطت إلى نحو 0.2\% وتكون
المقاومة قد ارتفعت إلى نحو 1\%. وفي اليوم السادس يعود المجموع المتقطع
إلى 100\%، مؤلفًا كله تقريبًا من \omterm{def:g10:cells-common-unit:cell}{خلايا} مقاومة (فالناجيات الحساسات من
اليوم الثالث نمت هي أيضًا، لكن من قاعدة أصغر).
\end{solution}

\begin{solution}{exo:g11:antibiotic-resistance:8}
\omterm{def:g10:universal-dna:gene}{المورّثة} الصبغية لا تنتقل إلا إلى ذرية النسيلة؛ أما البلازميد فينتقل
أفقيًا كذلك، إلى \omterm{def:g10:cells-common-unit:cell}{خلايا} لا قرابة بينها وإلى \omterm{def:g10:biodiversity-scales:species}{أنواع} أخرى، ونقلة واحدة
تعطي \omterm{def:g10:cells-common-unit:cell}{الخلية} مقاومة في الحال من دون انتظار \omterm{def:g11:mutations:mutation}{طفرة}.
\end{solution}

\begin{solution}{exo:g11:antibiotic-resistance:9}
في اليوم الثالث تكون \omterm{def:g10:cells-common-unit:cell}{الخلايا} الحساسة قد هبطت إلى كسر من المئة بينما
ترتفع المقاومة؛ ومن دون الدواء تعود الناجيات كلها إلى النمو، والجمهرة
العائدة أغنى \omterm{def:g10:cells-common-unit:cell}{بالخلايا} المقاومة من الأصلية بكثير. فالانتكاسة لا تعود
تستجيب للدواء نفسه.
\end{solution}

\begin{solution}{exo:g11:antibiotic-resistance:10}
تستعمل المستشفيات \omterm{def:g11:antibiotic-resistance:antibiotic}{المضادات الحيوية} استعمالًا مكثفًا على مرضى كثيرين
متقاربين: فالانتقاء مستمر والسلالات المنتقاة تنتقل من مريض إلى مريض
بالأيدي والأسطح.
\end{solution}

\begin{solution}{exo:g11:antibiotic-resistance:11}
يعمل الستربتوميسين بارتباطه \omterm{def:g11:gene-expression:protein}{ببروتين} ريبوزومي؛ \omterm{def:g11:gene-expression:protein}{والبروتين} المتغير لا يعود
يرتبط به فتستمر \omterm{prop:g11:gene-expression:translation}{الترجمة}. لكن التغير يضعف كذلك عمل \omterm{def:g10:cells-common-unit:organelle}{الريبوزوم} الطبيعي
قليلًا، فيصير تركيب \omterm{def:g11:gene-expression:protein}{البروتين}، ومنه النمو، أبطأ قليلًا --- وهي كلفة
المقاومة.
\end{solution}

\begin{solution}{exo:g11:antibiotic-resistance:12}
$10^{-8} \times 10^{-8} = 10^{-16}$ لكل انقسام: فمع $10^{12}$
بكتيريا، لا يرجح أن تحمل \omterm{def:g10:cells-common-unit:cell}{خلية} كليهما. وفي السل يدوم العلاج أشهرًا
والبكتيريا بالمليارات، فانتقاء المقاومة لدواء واحد مؤكد؛ أما دواءان
معًا فلا يتركان ناجية تنتقى.
\end{solution}

\begin{solution}{exo:g11:antibiotic-resistance:13}
الجرعات المنخفضة تنتقي سلالات مقاومة في معي الحيوانات من دون أن تقتل
الجمهرة؛ ويكتسبها العمال بالتماس؛ وتحمل البلازميدات \omterm{def:g10:universal-dna:gene}{المورّثات} إلى
الممرضات البشرية وإلى البيئة عبر السماد. فترتفع المقاومة للدواء،
وللأدوية البشرية القريبة منه، في الطب.
\end{solution}

\begin{solution}{exo:g11:antibiotic-resistance:14}
كثيرًا ما تكلف المقاومة البكتيريا شيئًا (ريبوزومًا أبطأ، أو إنزيمًا
يصنع): فمن دون الدواء تنمو \omterm{def:g10:cells-common-unit:cell}{الخلايا} الحساسة أسرع قليلًا وتغلب المقاومة
عبر أجيال كثيرة. فتهبط النسبة --- لكن \omterm{def:g10:universal-dna:gene}{المورّثات} المقاومة تبقى في
الجمهرة بتواتر منخفض، مستعدة لأن تنتقى من جديد.
\end{solution}

\begin{solution}{exo:g11:antibiotic-resistance:15}
البكتيريا لا تتكيف فرادى. ففي أي جمهرة كبيرة تحمل بضع \omterm{def:g10:cells-common-unit:cell}{خلايا}، \omterm{def:g11:mutations:mutation}{بطفرة}
عشوائية أو ببلازميد تلقته من \omterm{def:g10:cells-common-unit:cell}{خلية} أخرى، أليلًا يجعلها مقاومة. ويقتل
\omterm{def:g11:antibiotic-resistance:antibiotic}{المضاد الحيوي} غيرها؛ فتتكاثر \omterm{def:g10:cells-common-unit:cell}{الخلايا} المقاومة وتنقل \omterm{def:g10:universal-dna:gene}{الأليل} إلى ذريتها.
فالجمهرة تصير مقاومة لأنها فرزت، لا لأن أفرادها تعلموا شيئًا.
\end{solution}

\begin{solution}{pb:g11:antibiotic-resistance:1}
\textbf{1.} نحو $10^{10}$ انقسام (واحد لكل \omterm{def:g10:cells-common-unit:cell}{خلية} أنتجت).

\textbf{2.} $10^{10} \times 10^{-8} = 100$ \omterm{def:g10:cells-common-unit:cell}{خلية} مقاومة، وسطيًا.

\textbf{3.} مئة \omterm{def:g10:cells-common-unit:cell}{خلية} بين عشرة مليارات غير مرئية على طبق: فالمرج يسلك
سلوك الحساسة لأن 99.999999\% منه كذلك.

\textbf{4.} الأموكسيسيلين: حساسة (24 فوق 17). والإريثروميسين: مقاومة
(12 دون 20). والبنسلين: مقاومة تمامًا، ولا تثبيط.

\textbf{5.} البنسلين عديم الفائدة ضد هذه السلالة. والاختيار بين
الأدوية الحساسة يزن كذلك الآثار الجانبية والطيف؛ ويفضَّل أضيق دواء
فعال ليعفي سائر بكتيريا المريض.

\textbf{6.} أربع فترات من 6 ساعات في اليوم: $10^{10} \times 0.1^4 =
10^6$ بعد يوم؛ و$10^{10} \times 0.1^{12} = 0.01$ بعد 3 أيام --- أي لا
شيء في الجوهر.

\textbf{7.} $100 \times 2^{48} \approx \num{2.8e16}$: وهو عبث، أكثر مما
يتسع له الجرح. فالنمو يحده الحيز والمغذيات والجهاز المناعي.

\textbf{8.} من 100 \omterm{def:g10:cells-common-unit:cell}{خلية} إلى $10^{10}$ يلزم $\log_2 10^8 \approx 27$
مضاعفة، أي نحو 13 ساعة؛ وعمليًا يتحرر الحيز على مدى اليوم الأول، فتكون
النسيلة المقاومة قد حلت محل الجمهرة الحساسة بحلول اليوم الأول أو
الثاني.

\textbf{9.} الحساسة: تهبط بعامل 10 كل 6 ساعات حتى تنعدم بحلول اليوم
الثالث. والمقاومة: ترتفع من 100 لتملأ الحيز في يوم أو يومين، ثم تستوي
عند $10^{10}$. والمجموع: انخفاض أثناء اليوم الأول، ثم عودة إلى
$10^{10}$، وهو الآن مقاوم كله.

\textbf{10.} إذا صفّى الجهاز المناعي \omterm{def:g10:cells-common-unit:cell}{الخلايا} أسرع مما تستطيع المئة
المقاومة التكاثر، انتهت العدوى رغمها؛ وإذا فاق نمو النسيلة التصفية،
صارت العدوى مقاومة. \omterm{def:g11:antibiotic-resistance:antibiotic}{فالمضادات الحيوية} تعمل مع الجهاز المناعي، لا بدلًا
منه.

\textbf{11.} الحساسة بعد يومين: $10^{10} \times 0.1^8 = 100$. والمقاومة:
نحو $10^{10}$ إن كانت النسيلة قد ملأت الحيز، أو آلاف كثيرة على الأقل:
فالناجيات مقاومة كلها تقريبًا.

\textbf{12.} النمو بالنسبة نفسها يحفظ النسبة: أي مقاومة 100\% في
الجوهر.

\textbf{13.} لا أثر: فالنسيلة المقاومة تتجاهل الأموكسيسيلين.

\textbf{14.} الطبق الأول حمل الجمهرة الأصلية، و99.999999\% منها حساسة؛
أما طبق الانتكاسة فيحمل النسيلة المنتقاة، وهي مقاومة كلها، فينمو المرج
حتى القرص.

\textbf{15.} تستطيع \omterm{def:g10:universal-dna:gene}{المورّثة} الانتقال إلى \omterm{def:g10:biodiversity-scales:species}{أنواع} أخرى يحملها مرضى آخرون
أو عاملون، فتصير سائر ممرضات الجناح مقاومة من دون انتظار طفرتها
الخاصة.

\textbf{16.} في الجناح، 400 جرعة في السنة في جمهرة صغيرة مغلقة تفرز
البكتيريا باستمرار، والسلالات المنتقاة تنتقل بين المرضى؛ أما في المدينة
فالانتقاء مخفف بين ناس كثيرين \omterm{def:g10:biodiversity-scales:species}{وأنواع} كثيرة.

\textbf{17.} مخططات المضادات: اختيار دواء لا تستطيع السلالة مقاومته،
فلا تنتقى نسيلة. والجرعات الكاملة: فلا جرعة مقطوعة تترك الناجيات
المقاومة تنمو من جديد. والنظافة: فلا انتقال للنسيلة المنتقاة بين
المرضى. والعزل: فلا انتشار أفقي ولا شاقولي من الحاملين.

\textbf{18.} انتقاء أقل أتاح للسلالات الحساسة، وهي تنمو أسرع قليلًا،
أن تستعيد أرضًا، وقطع انتقال المقاومة. ولم تعد إلى 5\% لأن السلالات
المقاومة وبلازميداتها تبقى بتواتر منخفض وتنتقى من جديد مع كل جرعة
باقية.

\textbf{19.} كل جمهرة جرثومية كبيرة تحوي \omterm{def:g10:cells-common-unit:cell}{خلايا} مقاومة للدواء الجديد
\omterm{def:g11:mutations:mutation}{بطفرة} مصادفة؛ واستعماله ينتقيها، وفي بضع سنين تشيع السلالة المقاومة،
كما حدث مع كل \omterm{def:g11:antibiotic-resistance:antibiotic}{مضاد حيوي} أدخل حتى الآن.

\textbf{20.} نحو مئة \omterm{def:g10:cells-common-unit:cell}{خلية} مقاومة بين عشرة مليارات قبل أي دواء؛ وفي
اليوم الثاني من العلاج ورثت الجرح؛ وقطع الجرعة باكرًا حوّل العدوى
القابلة للشفاء إلى عدوى مقاومة انتشرت.
\end{solution}
